% Generated from project outputs. Do not edit by hand.
\begin{table}[!htbp]
\centering
\begin{threeparttable}
\caption{Including nonlabor by subsample: 2016-2019}
\label{tab:descriptive-2016-2019-subsamples-y3}
\small
\setlength{\tabcolsep}{4pt}
\renewcommand{\arraystretch}{1.15}
\begin{tabularx}{\linewidth}{@{}Xrrrrr@{}}
\toprule
Sample & N & Mean t & SD t & Mean t+1 & SD t+1 \\
\midrule
General & 276 & 2,822.6 & 1,947.4 & 2,867.7 & 1,869.6 \\
Men & 276 & 3,133.9 & 2,350.7 & 3,156.9 & 2,241.0 \\
Women & 276 & 2,289.0 & 1,527.4 & 2,383.1 & 1,569.6 \\
Urban & 276 & 3,103.9 & 1,990.7 & 3,131.6 & 1,889.2 \\
Rural & 271 & 2,108.7 & 1,478.6 & 2,248.4 & 1,609.0 \\
Indigenous & 271 & 2,083.8 & 1,425.3 & 2,144.7 & 1,400.6 \\
Non-Indigenous & 276 & 3,049.3 & 1,993.5 & 3,110.3 & 1,929.8 \\
Bottom 99 percent & 276 & 2,574.8 & 1,445.5 & 2,628.5 & 1,408.8 \\
Bottom 90 percent & 276 & 2,134.2 & 1,027.0 & 2,177.3 & 988.3 \\
Bottom 50 percent & 272 & 1,193.1 & 531.3 & 1,247.5 & 533.0 \\
\bottomrule
\end{tabularx}
\begin{tablenotes}[flushleft]
\footnotesize
\item Monthly income in quetzales. Unweighted descriptive statistics of survey-weighted cohort means. N counts cohort-transition rows, not individuals or independent cohorts. SD is dispersion across cohort means, not the standard error of a mean. Both incomes must be observed. Subsamples overlap; bottom-income groups are cumulative.
\end{tablenotes}
\end{threeparttable}
\end{table}
